\section{Introduction}
\label{sec:introduction}

Oral control sessions are among the most visible accountability rituals in
parliamentary democracies. In the Spanish Congreso, registered questions to the
President of the Government set a public agenda: a named deputy asks, and the
Prime Minister must respond on the floor. Yet political response is not
guaranteed. Speakers may answer directly, answer only in part, or talk around
the demand while remaining topically related. How often first responses are
explicitly responsive, and how that pattern varies across institutional
contexts of questioning, is an empirical measurement problem before it is a
causal one.

This paper studies \emph{reply explicitness} in Spanish prime-ministerial
control sessions. The unit of analysis is the exchange formed by a registered
oral question, the questioner's opening turn, and the Prime Minister's first
response. The outcome measure is a three-way human judgement of whether that
first response is an explicit reply, a partial reply, or an absence of reply,
anchored by the oral question and informed by the registered wording. The study
follows the reply / non-reply tradition associated with political-interview and
PMQs research \citep{bull_mayer_1993,bull_1994_identifying,bull_strawson_2020},
adapted to the Spanish plenary control format.

A fixed institutional baseline---the questioner's individual vote on the XIV
investiture that formed the government---is used only as a descriptive
grouping (\texttt{formation\_vote\_alignment}; primary contrast opposed versus
non-opposed). The design does not claim that alignment \emph{causes} reply
behaviour. Question-side features (form, topic, confrontational formulation)
are recorded independently of replies because they are strongly imbalanced by
alignment in this population; alignment differences therefore cannot be read as
clean causal effects of political side alone.

The contribution is a transparent measurement pipeline for Spanish control
sessions under pinned sources (ParlaMint-ES and Congreso initiative records),
an out-of-pool development and calibration sequence for the coding instrument,
and a pre-specified descriptive analysis of first-response explicitness. No
primacy claim is made relative to prior Spanish parliamentary pragmatics or UK
PMQs work.
